\documentclass[11pt]{article}
\usepackage[T1]{fontenc}
\usepackage{lmodern}
\usepackage{amsmath,amssymb,amsthm,mathtools}
\usepackage{enumitem}
\usepackage{graphicx}
\usepackage[margin=1in]{geometry}
\usepackage[colorlinks=true,linkcolor=blue,citecolor=blue,urlcolor=blue]{hyperref}
\setlength{\emergencystretch}{1em}

\newtheorem{theorem}{Theorem}
\newtheorem{lemma}{Lemma}
\newtheorem{corollary}{Corollary}
\newtheorem{remark}{Remark}
\newcommand{\R}{\mathbb{R}}
\newcommand{\C}{\mathbb{C}}
\newcommand{\K}{\mathcal{K}}
\newcommand{\F}{\mathcal{F}}
\newcommand{\Pcal}{\mathcal{P}}
\newcommand{\dd}{\,d}

\title{Krylov Solvability for Scalar Friedrichs Transport Operators:\\
Stable Groups, Quasi-Analytic Data, and a Smooth Obstruction}
\author{Run \texttt{fa\_banach\_001}, agent \texttt{agent\_lane\_09}}
\date{9 August 2026}

\begin{document}
\maketitle

\section*{Status and scope}

Candidate coefficient-wide strong partial resolution, likely valid, pending
specialist review.  Question~3 of Caruso and Michelangeli asks for an
investigation of Krylov solvability for
\[
 -\frac{d}{dx}+1
 \quad\hbox{and}\quad
 -\frac{d}{dx}+c(x)
\]
under suitable conditions on the datum and coefficient~\cite{CMFriedrichs}.
This packet treats the full scalar coefficient class singled out there: for
the maximal closed realization and
\[
 c\in L^2_{\mathrm{loc}}(\R),\qquad
 0<\delta\leq \operatorname{Re}c(x)\leq M<\infty\quad\text{a.e.},
\]
it proves Krylov solvability for every quasi-analytic datum, gives an exact
weighted-derivative criterion for every smooth datum, and shows the positive
class is dense.  The constant-coefficient construction from the preceding
version supplies a Schwartz datum for which solvability fails, so mere
$A$-smoothness is insufficient.

Because the source question is an open-ended request to ``investigate'' rather
than a precise universal conjecture, the packet is conservatively retained as
\texttt{partial}.  Mathematically, however, the coefficient side is no longer
left open: the theorem covers every scalar coefficient satisfying the source's
bounded coercive-real-part hypothesis.  What is not supplied is a simpler
closed-form classification of all smooth data beyond the exact closure
criterion.

\section*{Source question}

The source question occurs in Section~4, PDF page~11.  The crop includes the
constant operator (4.4), the variable operator (4.5), and Question~3.

\begin{center}
\includegraphics[width=\linewidth]{figures/open_problem_crop.png}
\end{center}

The question asks: ``To investigate the Krylov-solvability of the inverse
problem $Af=g$ induced by (4.4) or (4.5) (or variants) under suitable
conditions on $g$ and $c$.''  The paper then records
$\widehat g\in C_c^\infty(\R)$ as a sufficient condition for (4.4).

\section*{Definitions}

For a closed operator $A$ on a Hilbert space $H$, put
\[
 C^\infty(A)=\bigcap_{n\geq0}\operatorname{dom}(A^n),\qquad
 \K(A,g)=\operatorname{span}\{g,Ag,A^2g,\ldots\}.
\]
The inverse problem $Af=g$ is \emph{Krylov-solvable} when it has a solution in
$\overline{\K(A,g)}$~\cite{CMUnbounded,CMNBounded}.  A nonzero
$g\in C^\infty(A)$ is \emph{quasi-analytic for $A$} when
\begin{equation}\label{eq:qa}
 \sum_{n=1}^\infty \|A^n g\|^{-1/n}=\infty.
\end{equation}
This is the standard quasi-analytic-vector condition appearing in the
Denjoy--Carleman/operator-theoretic correspondence~\cite{Chernoff}.

\section*{Proof intuition}

There are two changes of viewpoint.  First, a general stable-group lemma says
that if $-A$ generates a group whose positive orbit decays exponentially,
then the inverse is the orbit integral
$A^{-1}g=\int_0^\infty e^{-tA}g\,dt$.  For a vector satisfying
\eqref{eq:qa}, every scalar orbit coefficient is quasi-analytic.  Its entire
Taylor jet at zero is determined by the Krylov vectors.  Any vector
orthogonal to the Krylov closure therefore annihilates the whole orbit, and
hence annihilates its integral.

Second, the integrating factor $e^{-C}$, where
$C(x)=\int_0^x c(s)\,ds$, does not give a bounded similarity on ordinary
$L^2$.  It instead gives an exact \emph{unitary} equivalence to differentiation
on the weighted space $L^2(e^{2\operatorname{Re}C}dx)$.  There the group is
ordinary right translation, and coercivity of $\operatorname{Re}c$ is exactly
what makes positive translations decay.  This proves the coefficient-wide
positive theorem and also turns arbitrary-data solvability into an exact
weighted derivative-approximation question.

The positive condition is not cosmetic.  At $c=1$, Fourier transformation
turns the problem into polynomial approximation in a datum-dependent measure.
An affine lognormal moment construction produces a Schwartz datum with a
nonzero polynomial-orthogonal witness that still pairs nontrivially with the
resolvent.  Its power norms grow at least like $e^{n^2}$, placing it strictly
outside the quasi-analytic class.

\section*{A stable-group Krylov theorem}

\begin{theorem}[Stable group plus quasi-analyticity]\label{thm:group}
Let $A$ be closed and densely defined on a Hilbert space $H$.  Suppose $-A$
generates a strongly continuous group $T(t)$ and, for some $M_0,\delta>0$,
\begin{equation}\label{eq:stable}
 \|T(t)\|\leq M_0e^{-\delta t}\qquad(t\geq0).
\end{equation}
Then $A$ is boundedly invertible.  For every nonzero $A$-quasi-analytic vector
$g$, the unique solution of $Af=g$ is Krylov-solvable, and
\begin{equation}\label{eq:orbit-inverse}
 A^{-1}g=\int_0^\infty T(t)g\,dt
 \in\overline{\K(A,g)}.
\end{equation}
\end{theorem}

\begin{proof}
The Bochner integral $R=\int_0^\infty T(t)\,dt$ converges in operator norm by
\eqref{eq:stable}.  The standard generator identity
\[
 (-A)\int_0^N T(t)x\,dt=T(N)x-x
\]
and closedness, followed by $N\to\infty$, give $ARx=x$ for all $x\in H$.
For $x\in\operatorname{dom}A$, differentiating the orbit and integrating gives
$RAx=x$.  Hence $R=A^{-1}$.

Fix a quasi-analytic $g$, put
$M_g=\overline{\K(A,g)}$, and let $0\ne h\perp M_g$.  Define on $\R$
\[
 \phi_h(t)=\langle T(t)g,h\rangle.
\]
The smooth-vector theorem for groups gives
\[
 \phi_h^{(n)}(t)=\langle T(t)(-A)^ng,h\rangle.
\]
On every compact interval $I$ there is a constant $C_I$ such that
\begin{equation}\label{eq:derivative-bound}
 \sup_{t\in I}|\phi_h^{(n)}(t)|
 \leq C_I\|A^ng\|\,\|h\|\qquad(n\geq0).
\end{equation}
Multiplication of the derivative majorants by a fixed positive constant does
not change divergence of their inverse-root sum.  More explicitly, for
$M_n=C_I\|A^ng\|\|h\|$ and
$\beta_n=\inf_{k\geq n}M_k^{1/k}$, one has
$\beta_n^{-1}\geq M_n^{-1/n}$; hence Carleman's criterion
$\sum_n\beta_n^{-1}=\infty$ holds.  The Denjoy--Carleman theorem therefore
makes $\phi_h$ quasi-analytic on every compact interval~\cite{Chernoff}.  Since
\[
 \phi_h^{(n)}(0)=\langle(-A)^ng,h\rangle=0\qquad(n\geq0),
\]
quasi-analyticity gives $\phi_h\equiv0$ on $\R$.  Thus $T(t)g$ is orthogonal
to $M_g^\perp$ and belongs to $M_g$ for every $t$.  Closedness of $M_g$ and
the norm-convergent integral in \eqref{eq:orbit-inverse} imply
$A^{-1}g\in M_g$.
\end{proof}

\section*{The scalar variable-coefficient model}

Let $c\in L^2_{\mathrm{loc}}(\R)$ satisfy
\begin{equation}\label{eq:c-hyp}
 0<\delta\leq\operatorname{Re}c(x)\leq M<\infty\quad\text{a.e.}
\end{equation}
Set
\[
 C(x)=\int_0^x c(s)\,ds,\qquad
 \omega(x)=e^{2\operatorname{Re}C(x)},\qquad
 H_\omega=L^2(\R,\omega(x)\,dx).
\]
On $H_\omega$ define
\[
 Bq=-q',\qquad
 \operatorname{dom}B=\{q\in H_\omega:q'\in H_\omega\},
\]
where the derivative is distributional.  Multiplication
\[
 U:H_\omega\to L^2(\R),\qquad Uq=e^Cq,
\]
is unitary.  Define $A_c=UBU^{-1}$.

\begin{theorem}[Coefficient-wide resolution]\label{thm:coefficient}
Under \eqref{eq:c-hyp}, $A_c$ is the maximal closed realization of
$-d/dx+c(x)$ and is boundedly invertible.  If $c\in L^\infty(\R)$, then
$\operatorname{dom}A_c=H^1(\R)$.  Moreover:
\begin{enumerate}[label=\textup{(\roman*)}]
\item Every nonzero $g\in C^\infty(A_c)$ satisfying
\begin{equation}\label{eq:weighted-qa}
 \sum_{n=1}^\infty
 \left\|\left(\frac d{dx}\right)^n(e^{-C}g)
 \right\|_{L^2(\omega dx)}^{-1/n}=\infty
\end{equation}
has a Krylov solution.
\item For arbitrary $g\in C^\infty(A_c)$, put $q=e^{-C}g$.  The unique
solution is
\begin{equation}\label{eq:solution-formula}
 f(x)=e^{C(x)}\int_x^\infty e^{-C(y)}g(y)\,dy,
\end{equation}
and it is a Krylov solution if and only if
\begin{equation}\label{eq:exact-weighted-criterion}
 x\longmapsto\int_x^\infty q(y)\,dy
 \quad\text{belongs to}\quad
 \overline{\operatorname{span}\{q,q',q'',\ldots\}}^{,H_\omega}.
\end{equation}
\item The $A_c$-analytic vectors, hence also the vectors satisfying
\eqref{eq:weighted-qa}, form a dense subspace of $L^2(\R)$.
\end{enumerate}
\end{theorem}

\begin{proof}
The unitary identity $\|Uq\|_2=\|q\|_{H_\omega}$ is immediate.  Since
$C'=c$ locally in the weak sense,
\[
 UBU^{-1}f=e^C\bigl(-(e^{-C}f)'\bigr)=-f'+cf.
\]
Thus the unitary definition gives precisely the maximal distributional
realization of the displayed expression.  When $c$ is bounded, multiplication
by $c$ is bounded on $L^2$, so its domain is $H^1(\R)$.

The group generated by $-B$ is right translation,
\[
 S(t)q(x)=q(x+t).
\]
Strong continuity follows first on compactly supported smooth functions and
then by density, using the local comparability of translates of $\omega$.
For $t\geq0$, the change of variable $y=x+t$ gives
\begin{align}
 \|S(t)q\|_{H_\omega}^2
 &=\int_\R |q(y)|^2\omega(y-t)\,dy\notag\\
 &=\int_\R |q(y)|^2\omega(y)
   \exp\left(-2\int_{y-t}^{y}\operatorname{Re}c(s)\,ds\right)\,dy
 \leq e^{-2\delta t}\|q\|_{H_\omega}^2.\label{eq:translation-decay}
\end{align}
For negative time the analogous estimate is
$\|S(t)\|\leq e^{M|t|}$, so $S$ is a $C_0$ group.  Conjugating by $U$ gives
\begin{equation}\label{eq:physical-group}
 T_c(t)g(x)=
 \exp\left(-\int_0^t c(x+s)\,ds\right)g(x+t),
\end{equation}
and \eqref{eq:translation-decay} gives
$\|T_c(t)\|\leq e^{-\delta t}$ for $t\geq0$.  Theorem~\ref{thm:group}
applies.

If $g\in C^\infty(A_c)$ and $q=U^{-1}g=e^{-C}g$, unitary equivalence gives
\begin{equation}\label{eq:powers}
 \|A_c^ng\|_2=\|B^nq\|_{H_\omega}=\|q^{(n)}\|_{H_\omega},
 \qquad
 U^{-1}\overline{\K(A_c,g)}
 =\overline{\operatorname{span}\{q,q',q'',\ldots\}}^{,H_\omega}.
\end{equation}
Thus \eqref{eq:weighted-qa} is exactly \eqref{eq:qa}, and part~(i) follows.
Also,
\[
 B^{-1}q=\int_0^\infty S(t)q\,dt,
\]
which has the almost-everywhere representative
$B^{-1}q(x)=\int_x^\infty q(y)\,dy$.  This integral is absolutely convergent
for each fixed $x$ by Cauchy--Schwarz and the exponential lower growth of
$\omega$.  Applying $U$ gives \eqref{eq:solution-formula}; comparing it with
\eqref{eq:powers} proves the exact criterion in part~(ii).

For completeness, standard Gaussian group regularization supplies density in
part~(iii).  If $x\in H$ and $\rho_s$ is a centered Gaussian approximate
identity, then
\[
 x_s=\int_\R \rho_s(t)T_c(t)x\,dt\longrightarrow x.
\]
Integration by parts expresses $A_c^nx_s$ using $\rho_s^{(n)}$; Gaussian
derivative estimates give
$\|A_c^nx_s\|\leq C_sD_s^n\sqrt{n!}\,\|x\|$.  Hence $x_s$ is an analytic
vector.  Analytic vectors are quasi-analytic because Stirling's formula makes
the inverse-root series dominate a constant multiple of the harmonic series.
This is the familiar analytic-vector regularization mechanism~\cite{Nelson}.
\end{proof}

\section*{Sharpness at the constant coefficient}

The next theorem records the full constant-model classification and the
smooth obstruction from the preceding packet version.  Let $\F$ be the
unitary Fourier transform with $\F(f')=i\xi\widehat f$.  For
$A=-d/dx+1$ and nonzero $g\in C^\infty(A)$, put
\[
 d\mu_g(\xi)=|\widehat g(\xi)|^2\,d\xi,
 \qquad r(\xi)=\frac1{1-i\xi},
\]
and write $\Pcal$ for the complex polynomials.

\begin{theorem}[Exact constant criterion and obstruction]\label{thm:constant}
For $A=-d/dx+1$ on $H^1(\R)\subset L^2(\R)$:
\begin{enumerate}[label=\textup{(\roman*)}]
\item The solution of $Af=g$ is Krylov if and only if
$r\in\overline{\Pcal}^{L^2(\mu_g)}$.
\item If $\int e^{\varepsilon|\xi|}\,d\mu_g(\xi)<\infty$ for some
$\varepsilon>0$, then polynomials are dense in $L^2(\mu_g)$ and the problem
is Krylov-solvable.  In particular, bounded Fourier support is sufficient
without any smoothness assumption on $\widehat g$.
\item There is a nonzero $g\in\mathcal S(\R)$ for which the solution is not
Krylov.  This datum is not quasi-analytic for $A$.
\end{enumerate}
\end{theorem}

\begin{proof}
Under $\F$, $A$ is multiplication by $m(\xi)=1-i\xi$ and $A^{-1}$ is
multiplication by $r=m^{-1}$.  The map
\[
 J_g:L^2(\mu_g)\to L^2(\R),\qquad J_gp=p\widehat g,
\]
is an isometry with closed range.  Since polynomials in $m$ and in $\xi$
coincide,
\[
 \F\bigl(\overline{\K(A,g)}\bigr)
 =J_g\left(\overline{\Pcal}^{L^2(\mu_g)}\right),
 \qquad \widehat{A^{-1}g}=J_gr.
\]
This proves~(i).

For~(ii), if $h\in L^2(\mu_g)$ is polynomial-orthogonal, the finite complex
measure $d\eta=\overline h\,d\mu_g$ has an exponential moment of half the
given order by Cauchy--Schwarz.  Its bilateral Laplace transform is holomorphic
on a strip and every derivative at zero vanishes.  The identity theorem and
uniqueness of Fourier transforms of finite measures give $\eta=0$
\cite[Chapter~9]{RudinRCA}.  Thus the polynomial orthogonal complement is
trivial.

For~(iii), let $\nu$ be the lognormal probability measure
\begin{equation}\label{eq:lognormal}
 d\nu(x)=\frac1{x\sqrt{2\pi}}
 \exp\left(-\frac{(\log x)^2}{2}\right)dx\qquad(x>0)
\end{equation}
and $h_0(x)=\sin(2\pi\log x)$.  If $X=e^Y$ with $Y$ standard normal, then
for every integer $n\geq0$,
\begin{equation}\label{eq:zero-moments}
 \int_0^\infty x^nh_0(x)\,d\nu(x)
 =\operatorname{Im}\exp\left(\frac{(n+2\pi i)^2}{2}\right)=0.
\end{equation}
The signed measure $\sigma=h_0\nu$ is nonzero, so its Cauchy transform
\[
 F(z)=\int_0^\infty\frac{1}{x-z}\,d\sigma(x),\qquad\operatorname{Im}z<0,
\]
is not identically zero.  Indeed, otherwise its Poisson convolutions would
all vanish.  The approximate-identity limit would then force $\sigma=0$.
Choose $z_0=u+iv$ with $v<0$ and $F(z_0)\ne0$.  Put
\[
 a=-1/v>0,\qquad b=-au,\qquad T(x)=ax+b,
\]
and transport $\nu$ by $T$.  If $\mu=T_\#\nu$ and $h(Tx)=h_0(x)$,
equation \eqref{eq:zero-moments} says that $h$ annihilates every polynomial in
$L^2(\mu)$, whereas
\[
 \int r(\xi)h(\xi)\,d\mu(\xi)
 =\frac{i}{a}F(z_0)\ne0,
\]
because $1-i(ax+b)=-ia(x-z_0)$.  Hence
$r\notin\overline{\Pcal}^{L^2(\mu)}$.

Let $w$ be the density of $\mu$ and choose $\widehat g=\sqrt w$.  Before the
affine change of variable, the square-root density, extended by zero to the
negative half-line, is a constant multiple of
\[
 q_0(x)=x^{-1/2}\exp\left(-\frac{(\log x)^2}{4}\right)\qquad(x>0).
\]
Every derivative has the form
$x^{-k-1/2}P_k(\log x)e^{-(\log x)^2/4}$.  The log-Gaussian dominates every
exponential as $\log x\to\pm\infty$, proving flatness at zero and Schwartz
decay at infinity.  Therefore $\widehat g$, and hence $g$, is Schwartz and
$\mu_g=\mu$.  Part~(i) now proves non-solvability.

Finally this example is outside the positive class of
Theorem~\ref{thm:coefficient}.  With $\xi=aX+b$, for all large $n$ and on
$Y=\log X\in[2n,2n+1]$, one has $|\xi|\geq aX/2$ and
\[
 \int_{2n}^{2n+1}e^{2ny-y^2/2}\frac{dy}{\sqrt{2\pi}}
 \geq c_0e^{2n^2}.
\]
Consequently
$\|A^ng\|\geq c_1(a/2)^ne^{n^2}$ for all large $n$, so
$\sum_n\|A^ng\|^{-1/n}<\infty$.  The Schwartz datum is $A$-smooth but not
$A$-quasi-analytic.
\end{proof}

\section*{Verification notes}

The following points received separate adversarial checks.
\begin{itemize}
\item The integrating factor is used as a unitary map from the weighted
space, not incorrectly as a bounded similarity on unweighted $L^2$.
\item The translation estimate has the correct orientation: right
translation compares $\omega(y-t)$ with $\omega(y)$ and decays for positive
time because $\operatorname{Re}c\geq\delta$.
\item The Denjoy--Carleman argument is applied to scalar group orbits on every
compact real interval; the group, rather than only a semigroup, removes an
endpoint uniqueness issue.
\item The maximal-domain statement follows by unitary definition and the
weak product rule.  The familiar $H^1$ domain is asserted only when $c$ is
bounded.
\item The negative example exhibits both a polynomial-orthogonal witness and
a nonzero pairing with the specific resolvent; moment indeterminacy alone is
not used.
\item The lower power-norm bound confirms that the counterexample does not
contradict the quasi-analytic theorem.
\end{itemize}

The accompanying script checks the lognormal moment formula at high precision
and evaluates one sample Cauchy pairing as a regression test.  These numerical
checks are not used in the proof.  The detailed report is
\texttt{verification.md}.

\section*{Novelty check and limitations}

The run indexes and a bounded arXiv/web search were checked using the exact
source title and authors, ``Krylov solvability Friedrichs'',
``$-d/dx+c(x)$ Krylov'', ``analytic vector Krylov-solvable'', stable
semigroups, and quasi-analytic vectors.  The closest works inspected were
arXiv:1811.08202, 2001.08127, 2102.13626, and 2210.04752
\cite{CMNBounded,CMUnbounded,CMPerturb,CarusoNormal}.  They contain the
general framework, self-adjoint/skew-adjoint mechanisms, perturbation theory,
and bounded or compact normal examples.  No later paper explicitly resolving
Question~3, and no statement of Theorem~\ref{thm:group} in this Krylov form,
was located in this bounded search.  This is not proof of novelty.

The stable-group theorem is a short application of the classical
Denjoy--Carleman theorem and may be folklore.  The candidate-new contribution
is its coefficient-wide Friedrichs application together with the exact
weighted model and the quasi-analytic-versus-smooth separation.  The
coefficient may have unbounded imaginary part under the stated local
regularity, but the realization is then the maximal domain rather than
necessarily $H^1$.  Matrix-valued Friedrichs systems and a simpler intrinsic
classification of every smooth datum remain outside the packet.

\section*{Human-review recommendation}

Preserve as a coefficient-wide strong partial-resolution candidate.  A
reviewer should focus on the precise Denjoy--Carleman criterion used after
\eqref{eq:derivative-bound}, the maximal-domain identification under only
$c\in L^2_{\mathrm{loc}}$, and the density estimate in Gaussian group
regularization.  If the natural interpretation of the open-ended source
question is ``give useful sharp conditions across the scalar class'', this
packet can reasonably be viewed as a full scalar resolution.  Protocol-wise,
\texttt{partial} remains the safer label pending that judgment.

\begin{thebibliography}{9}

\bibitem{CMFriedrichs}
N.~A. Caruso and A. Michelangeli,
\emph{Open problems and perspectives on solving Friedrichs systems by Krylov
approximation}, arXiv:2406.05777v2 (2024).

\bibitem{CMUnbounded}
N.~A. Caruso and A. Michelangeli,
\emph{Krylov solvability of unbounded inverse linear problems},
Integral Equations Operator Theory 93 (2021), Paper No.~1;
arXiv:2001.08127.

\bibitem{CMNBounded}
N.~A. Caruso, A. Michelangeli, and P. Novati,
\emph{On Krylov solutions to infinite-dimensional inverse linear problems},
Calcolo 56 (2019), Paper No.~32; arXiv:1811.08202.

\bibitem{CMPerturb}
N.~A. Caruso and A. Michelangeli,
\emph{Krylov solvability under perturbations of abstract inverse linear
problems}, Journal of Applied Analysis 29 (2023), 3--29;
arXiv:2102.13626.

\bibitem{CarusoNormal}
N.~A. Caruso,
\emph{A note on the Krylov solvability of compact normal operators on Hilbert
space}, Complex Analysis and Operator Theory 17 (2023);
arXiv:2210.04752.

\bibitem{Chernoff}
P.~R. Chernoff,
\emph{Quasi-analytic vectors and quasi-analytic functions},
Bulletin of the American Mathematical Society 81 (1975), 637--646.

\bibitem{Nelson}
E.~Nelson,
\emph{Analytic vectors}, Annals of Mathematics (2) 70 (1959), 572--615.

\bibitem{RudinRCA}
W.~Rudin, \emph{Real and Complex Analysis}, 3rd ed., McGraw--Hill, 1987.

\end{thebibliography}

\end{document}
